% =============================================================================
%  Plan de pacing — Nice by UTMB 100M — module Locomotion Twin
%  Compiler : latexmk -xelatex main.tex   (XeLaTeX + biber)
% =============================================================================
\documentclass{locomotionreport}

\renewcommand{\LLCoverLogoPath}{assets/logo_primary_deep}
\renewcommand{\LLHeaderLogoPath}{assets/logo_full_white}

% *Math — macros reprises du thème Beamer Locomotion Lab.
% (amsmath, amssymb, mathtools, bm, stmaryrd sont déjà chargés par la classe.)
\usepackage{nicematrix}
\usepackage{indentfirst}

%\newcommand{\vect}[1]{\text{\bfseries {#1}}}
\newcommand{\vect}[1]{\bm{#1}}
\renewcommand{\d}{\mathrm{d}}
\newcommand{\vmax}{V_{\text{max}}}
\newcommand{\imid}{_{k+\frac{1}{2}}}
\newcommand{\immid}{_{k-\frac{1}{2}}}
\newcommand{\ione}{_{k+1}}
\newcommand{\imone}{_{k-1}}
\newcommand{\none}{^{n+1}}
\newcommand{\odx}[1]{\frac{\partial #1}{\partial x}}
\newcommand{\odxx}[1]{\frac{\partial^2 #1}{\partial x^2}}
\newcommand{\ody}[1]{\frac{\partial #1}{\partial y}}
\newcommand{\odz}[1]{\frac{\partial #1}{\partial z}}
\newcommand{\odt}[1]{\frac{\partial #1}{\partial t}}
\newcommand{\dt}[1]{\frac{\mathrm{d} #1}{\mathrm{d}t}}
\newcommand{\ok}{{\text{K}}}
\newcommand{\inter}{\llbracket 1, N \rrbracket}
\newcommand{\interr}{\llbracket 1, 2N \rrbracket}

\newcommand{\upo}{\bm{\upsilon}^o}
\newcommand{\upe}{\bm{\upsilon}^e}
\newcommand{\upor}{\bm{\upsilon}^o_\Re}
\newcommand{\upoi}{\bm{\upsilon}^o_\Im}
\newcommand{\uper}{\bm{\upsilon}^e_\Re}
\newcommand{\upei}{\bm{\upsilon}^e_\Im}


\newcommand{\A}{\mathcal{A}}
\newcommand{\B}{\mathcal{B}}
\newcommand{\C}{\mathcal{C}}
\newcommand{\Ckk}{\mathcal{C}^{\text{K}}}
\newcommand{\ck}{c^{\text{K}}}
\newcommand{\Ckkk}{\mathcal{C}^{\text{K1}}}

\newcommand{\D}{\mathcal{D}}
\newcommand{\E}{\mathcal{E}}
\newcommand{\F}{\mathcal{F}}
\newcommand{\G}{\mathcal{G}}
\newcommand{\HH}{\mathcal{H}}
\newcommand{\II}{\mathcal{I}}
\newcommand{\J}{\mathcal{J}}
\newcommand{\K}{\mathcal{K}}
\newcommand{\LL}{\mathcal{L}}
\newcommand{\M}{\mathcal{M}}
\newcommand{\N}{\mathcal{N}}
\newcommand{\OO}{\mathcal{O}}
\newcommand{\PP}{\mathcal{P}}
\newcommand{\QQ}{\mathcal{Q}}
\newcommand{\R}{\mathcal{R}}
\newcommand{\SSS}{\mathcal{S}}
\newcommand{\T}{\mathcal{T}}
\newcommand{\U}{\mathcal{U}}
\newcommand{\V}{\mathcal{V}}
\newcommand{\W}{\mathcal{W}}
\newcommand{\X}{\mathcal{X}}
\newcommand{\Y}{\mathcal{Y}}
\newcommand{\Z}{\mathcal{Z}}

\newcommand{\Gv}{\bm{G}}
\newcommand{\Rv}{\bm{R}}
\newcommand{\bv}{\bm{b}}
\newcommand{\q}{\bm{q}}
\newcommand{\qb}{\bm{q}_b}
\newcommand{\qp}{\bm{q}'}
\newcommand{\qh}{\widehat{\bm{q}}}
\newcommand{\qt}{\widetilde{\bm{q}}}
\newcommand{\Q}{\bm{Q}}
\newcommand{\Qb}{\bm{Q}_b}
\newcommand{\Qp}{\bm{Q}'}
\newcommand{\Qh}{\widehat{\bm{Q}}}
\newcommand{\Qt}{\widetilde{\bm{Q}}}
\newcommand{\Wv}{\bm{W}}
\newcommand{\Wb}{\bm{W}_b}
\newcommand{\Wp}{\bm{W}'}
\newcommand{\Wh}{\widehat{\bm{W}}}
%\newcommand{\Wt}{\widetilde{\bm{W}}}
\newcommand{\Wt}{\widehat{\bm{W}}^\dagger}
%\newcommand{\f}{\bm{F}}
%\newcommand{\fh}{\bm{\widehat{F}}}
\newcommand{\f}{\bm{\Phi}}
\newcommand{\fh}{\bm{\widehat{\Phi}}}
\newcommand{\g}{\bm{g}}
\newcommand{\uu}{\bm{u}}
\newcommand{\uh}{\bm{\widehat{u}}}
\newcommand{\rv}{\bm{r}}
\newcommand{\Xv}{\bm{X}}
\newcommand{\xv}{\bm{x}}
\newcommand{\xvr}{\bm{x}_{\Re}}
\newcommand{\xvi}{\bm{x}_{\Im}}
\newcommand{\vv}{\bm{v}}
\newcommand{\vvr}{\bm{v}_{\Re}}
\newcommand{\vvi}{\bm{v}_{\Im}}
\newcommand{\ww}{\bm{w}}
\newcommand{\wh}{\widehat{\bm{w}}}
\newcommand{\wpp}{\bm{w}'}
\newcommand{\vvt}{\bm{\widetilde{v}}}
\newcommand{\wwt}{\bm{\widetilde{w}}}
\newcommand{\yv}{\bm{y}}
\newcommand{\zv}{\bm{z}}
\newcommand{\uodd}{\bm{u}^o}
\newcommand{\ueven}{\bm{u}^e}
\newcommand{\hh}{\hat{\bm{h}}}
\newcommand{\bh}{\hat{\bm{b}}}
\newcommand{\oof}{\text{F}}
\newcommand{\oob}{\text{B}}

\newcommand{\snl}{\sigma^{\text{NL}}}
\newcommand{\stnl}{\text{St}^{\text{NL}}}

\newcommand{\CFL}{\text{CFL}}
\newcommand{\RE}{\text{Re}}
\newcommand{\PR}{\text{Pr}}
\newcommand{\MA}{\text{M}}
\newcommand{\ST}{\text{ST}}

\newcommand{\Sp}{\text{Sp}}

\newcommand*{\tensor}[1]{\overline{\overline{#1}}}

%\floatname{algorithm}{Algorithme}
% Adapté : ne redéfinir que si le paquet algorithmic/algorithmicx est chargé.
\ifdefined\algorithmicrequire
  \renewcommand{\algorithmicrequire}{\textbf{Input:}}
  \renewcommand{\algorithmicensure}{\textbf{Output:}}
\fi

\addbibresource{references.bib}

% --- macros légères propres à ce rapport ---
\newcommand{\Deq}{D_{\mathrm{eq}}}
\newcommand{\kmh}{\,km\,h\textsuperscript{-1}}
\newcommand{\VC}{\mathrm{VC}}

\title{Plan de pacing personnalis\'e}
\subtitle{Nice C\^ote d'Azur by UTMB\textsuperscript{\textregistered} \textperiodcentered{} 100M \textperiodcentered{} 165\,km \textperiodcentered{} 8\,900\,m\,D+ \textperiodcentered{} 25 septembre 2026}
\author{The Locomotion Lab}
\date{24 juin 2026}
\reporttype{Rapport d'analyse \textperiodcentered{} prototype}
\reportref{LL-TWIN-2026-001}
\reportversion{v0.9}
\reportclient{Valentin \textperiodcentered{} coureur}

\begin{document}

\makecover

\begin{llabstract}
Ce rapport construit un \textbf{plan d'allure individualis\'e} pour le 100M du Nice C\^ote d'Azur by UTMB\textsuperscript{\textregistered}
(165\,km, 8\,900\,m\,D+) \`a partir de deux sources : la trace GPX du parcours et 449~fichiers d'entra\^inement issus
de votre montre Coros (18~mois, du 6~janvier~2025 au 22~juin~2026). La cha\^ine \emph{Locomotion Twin} estime votre
profil physiologique --- vitesse critique, exposant d'endurance, durabilit\'e --- puis le confronte au co\^ut \'energ\'etique
de la pente le long du parcours pour pr\'edire un temps d'arriv\'ee et r\'epartir l'effort entre les seize ravitaillements.
La pr\'ediction centrale est de \textbf{30\,h\,24} (intervalle 80\,\% : 29\,h\,36 -- 31\,h\,18). La m\'ethode est valid\'ee
\emph{hors \'echantillon} sur vos propres ultras, avec une erreur moyenne de \textbf{2,8\,\%}.
\keywords{pacing \textperiodcentered{} ultra-trail \textperiodcentered{} vitesse critique \textperiodcentered{} durabilit\'e \textperiodcentered{} incertitude}
\end{llabstract}

\LLtoc
\clearpage

% =====================================================================
\section{Synth\`ese}

Ce document est un \emph{prototype} de rapport client produit par le module \emph{Locomotion Twin}. Il r\'epond \`a une
question simple --- \,\guillemotleft\,\`a quelle allure courir le 100M, kilom\`etre par kilom\`etre, pour aller au bout
dans les meilleures conditions\,\guillemotright\, --- en s'appuyant uniquement sur des donn\'ees mesur\'ees et des
mod\`eles physiologiques publi\'es. Aucune valeur n'est \og~au jug\'e~\fg : chaque chiffre d\'ecoule de vos fichiers
d'entra\^inement ou de la trace officielle.

\begin{llnote}[L'essentiel]
\textbf{Temps d'arriv\'ee pr\'edit : 30\,h\,24}, soit une arriv\'ee le \textbf{samedi vers 19\,h\,23} pour un d\'epart
le vendredi 25~septembre \`a 13\,h\,00. Intervalle de confiance \`a 80\,\% : \textbf{29\,h\,36 -- 31\,h\,18}.
La barri\`ere horaire (50\,h) laisse une marge consid\'erable.
\\[2pt]
\textbf{Intensit\'e cible : $\approx$\,63\,\% de votre vitesse critique}, soutenue avec une d\'erive ma\^itris\'ee
($-15\,\%$ entre le d\'ebut et la fin). \`A cette intensit\'e, le facteur limitant n'est pas la \og~r\'eserve
ana\'erobie\,\fg\ mais la \textbf{durabilit\'e} (l'usure progressive) et le \textbf{ravitaillement}.
\\[2pt]
\textbf{Une seule nuit} de course pour le sc\'enario m\'edian : frontale n\'ecessaire du \textbf{km\,38 (ven. 19\,h\,45)}
au \textbf{km\,94 (sam. 7\,h\,06)} environ.
\end{llnote}

Les trois piliers du raisonnement sont d\'evelopp\'es ci-dessous : (i)~le parcours et son co\^ut, (ii)~votre jumeau
num\'erique issu des donn\'ees, (iii)~la pr\'ediction, sa validation, puis le plan segment par segment. Les limites
sont \'enonc\'ees franchement en fin de rapport : un plan n'est utile que si l'on sait \emph{ce qu'il ne capture pas}.

% =====================================================================
\section{Le parcours et ses exigences}

\subsection{Donn\'ees officielles}
Le 100M relie \textbf{Auron} \`a la \textbf{Promenade des Anglais} \`a Nice. Les caract\'eristiques officielles sont
de 165\,km pour 8\,900\,m de d\'enivel\'e positif, d\'epart le vendredi 25~septembre~2026 \`a 13\,h\,00, barri\`ere
horaire de 50~heures, point culminant au lac de Rabuons \`a plus de 2\,650\,m~\cite{niceutmb}. C'est une course de
\og~haute montagne\,\fg\ en premi\`ere moiti\'e, puis une longue descente vers la mer.

\subsection{Le profil r\'eel issu de votre trace}
La trace GPX fournie compte 6\,077~points. Apr\`es lissage de l'altim\'etrie (fen\^etre glissante de 150\,m,
calibr\'ee pour reproduire le d\'enivel\'e officiel), on obtient un d\'enivel\'e de \textbf{8\,874\,m\,D+} et
\textbf{10\,461\,m\,D--} : la course \emph{descend} davantage qu'elle ne monte (d\'epart \`a 1\,591\,m, arriv\'ee
\`a 6\,m). La distance mesur\'ee en 3D (165,8\,km) co\"incide avec l'annonce officielle ; les rep\`eres kilom\'etriques
ci-apr\`es sont align\'es sur le carnet de route (arriv\'ee km\,167,2).

\begin{figure}[ht]
  \centering
  \includegraphics[width=\linewidth]{figures/profil.png}
  \caption{Profil altim\'etrique liss\'e. Points dor\'es : les seize ravitaillements. La grande bosse initiale
  (Auron \(\rightarrow\) Rabuons) concentre l'essentiel de la difficult\'e verticale ; la seconde moiti\'e encha\^ine
  les \og~quatre monts\,\fg\ avant la mer.}
  \label{fig:profil}
\end{figure}

\subsection{Le co\^ut de la pente : la distance \'equivalente}
Courir n'a pas le m\^eme co\^ut \`a plat et en c\^ote. On le quantifie par le \textbf{co\^ut \'energ\'etique de la
course} \(C_r\) en fonction de la pente \(i\), mesur\'e exp\'erimentalement par Minetti et~al.~\cite{minetti2002} :
\begin{equation}
C_r(i) = 155{,}4\,i^5 - 30{,}4\,i^4 - 43{,}3\,i^3 + 46{,}3\,i^2 + 19{,}5\,i + 3{,}6,
\label{eq:minetti}
\end{equation}
o\`u \(C_r\) est en J\,kg\textsuperscript{-1}\,m\textsuperscript{-1} et vaut 3,6 \`a plat. Le rapport \(f(i)=C_r(i)/3{,}6\)
indique \og~combien de m\`etres \`a plat co\^ute un m\`etre de pente \(i\)\,\fg. En int\'egrant \(f\) le long du parcours,
on obtient sa \textbf{distance \'equivalente \`a plat} :
\begin{equation}
\Deq = \int_0^{L} f\big(i(x)\big)\,\mathrm{d}x \;\approx\; \textbf{200,1\,km}.
\label{eq:deq}
\end{equation}
Autrement dit, ces 167\,km de montagne \og~co\^utent\,\fg\ comme \textbf{200\,km \`a plat}. C'est cette grandeur,
et non la distance brute, qui gouverne le temps de course.

\subsection{D\'ecoupage par ravitaillements}
Le tableau~\ref{tab:demande} d\'etaille la demande de chacun des seize segments. Deux ressortent nettement :
le \textbf{segment\,2} (km\,8,1\,\(\rightarrow\)\,16,5), montée de 1\,303\,m vers Rabuons \`a \(+15{,}9\,\%\) de moyenne ;
et le \textbf{segment\,4} (km\,28,9\,\(\rightarrow\)\,38), descente \og~sauvage et technique\,\fg\ de 1\,757\,m \`a
\(-18{,}9\,\%\). La figure~\ref{fig:demande} visualise cette r\'epartition.

\begin{table}[ht]
  \centering
  \footnotesize
  \setlength{\tabcolsep}{4.5pt}
  \LLzebra
  \renewcommand{\arraystretch}{1.32}
  \begin{tabular}{@{}clrrrrrr@{}}
    \LLheaderrow
    \LLthead{\#} & \LLthead{Ravitaillement} & \LLthead{dist.} & \LLthead{D+} & \LLthead{D--} & \LLthead{$\Deq$} & \LLthead{pente} & \LLthead{alt. fin} \\
     & & \footnotesize(km) & \footnotesize(m) & \footnotesize(m) & \footnotesize(km) & \footnotesize(\%) & \footnotesize(m) \\
1 & R1 \textperiodcentered{} km\,8,1 & \LLnum{8,1} & \LLnum{215} & \LLnum{656} & \LLnum{7,7} & \LLnum{-5,6} & \LLnum{1152} \\
2 & R2 \textperiodcentered{} km\,16,5 & \LLnum{8,4} & \LLnum{1303} & \LLnum{19} & \LLnum{18,0} & \LLnum{15,9} & \LLnum{2436} \\
3 & R3 \textperiodcentered{} km\,28,9 & \LLnum{12,4} & \LLnum{692} & \LLnum{608} & \LLnum{15,0} & \LLnum{0,7} & \LLnum{2521} \\
4 & R4 \textperiodcentered{} km\,38 & \LLnum{9,1} & \LLnum{108} & \LLnum{1757} & \LLnum{7,2} & \LLnum{-18,9} & \LLnum{872} \\
5 & R5 \textperiodcentered{} km\,50,1 & \LLnum{12,1} & \LLnum{612} & \LLnum{788} & \LLnum{15,2} & \LLnum{-1,3} & \LLnum{695} \\
6 & R6 \textperiodcentered{} km\,63,7 & \LLnum{13,6} & \LLnum{807} & \LLnum{953} & \LLnum{16,3} & \LLnum{-1,2} & \LLnum{549} \\
7 & R7 \textperiodcentered{} km\,68,4 & \LLnum{4,7} & \LLnum{453} & \LLnum{59} & \LLnum{7,6} & \LLnum{8,6} & \LLnum{942} \\
8 & R8 \textperiodcentered{} km\,83,5 & \LLnum{15,1} & \LLnum{1048} & \LLnum{748} & \LLnum{19,9} & \LLnum{2,0} & \LLnum{1242} \\
9 & R9 \textperiodcentered{} km\,93,8 & \LLnum{10,3} & \LLnum{868} & \LLnum{250} & \LLnum{16,1} & \LLnum{6,1} & \LLnum{1861} \\
10 & R10 \textperiodcentered{} km\,109,8 & \LLnum{16,0} & \LLnum{267} & \LLnum{1367} & \LLnum{13,6} & \LLnum{-7,1} & \LLnum{761} \\
11 & R11 \textperiodcentered{} km\,121,8 & \LLnum{12,0} & \LLnum{491} & \LLnum{682} & \LLnum{12,8} & \LLnum{-1,6} & \LLnum{570} \\
12 & R12 \textperiodcentered{} km\,128,3 & \LLnum{6,5} & \LLnum{855} & \LLnum{188} & \LLnum{12,3} & \LLnum{10,6} & \LLnum{1237} \\
13 & R13 \textperiodcentered{} km\,140,4 & \LLnum{12,1} & \LLnum{196} & \LLnum{1056} & \LLnum{9,7} & \LLnum{-7,2} & \LLnum{376} \\
14 & R14 \textperiodcentered{} km\,147,4 & \LLnum{7,0} & \LLnum{249} & \LLnum{539} & \LLnum{7,0} & \LLnum{-4,3} & \LLnum{87} \\
15 & R15 \textperiodcentered{} km\,156,5 & \LLnum{9,1} & \LLnum{537} & \LLnum{242} & \LLnum{11,8} & \LLnum{3,3} & \LLnum{382} \\
16 & Arriv\'ee \textperiodcentered{} km\,167,2 & \LLnum{10,7} & \LLnum{173} & \LLnum{549} & \LLnum{9,9} & \LLnum{-3,6} & \LLnum{6} \\
  \end{tabular}
  \caption{Demande de chaque segment inter-ravitaillement (g\'eom\'etrie issue de la trace). \(\Deq\) : distance
  \'equivalente \`a plat ; \og~pente\,\fg\ : pente moyenne sign\'ee.}
  \label{tab:demande}
\end{table}

\begin{figure}[ht]
  \centering
  \includegraphics[width=\linewidth]{figures/demande.png}
  \caption{Demande verticale (haut) et distance r\'eelle vs \'equivalent plat (bas) par segment. L'\'ecart entre
  les deux barres du bas mesure la \og~taxe\,\fg\ de la pente.}
  \label{fig:demande}
\end{figure}

% =====================================================================
\section{Votre jumeau num\'erique}

Le \emph{jumeau} est un r\'esum\'e chiffr\'e de vos capacit\'es, extrait de vos donn\'ees. Sur les 449~activit\'es
fournies, \textbf{419 courses \`a pied} exploitables ont \'et\'e retenues : 532~jours couverts, 4\,957\,km bruts
(5\,621\,km en \'equivalent plat), 741~heures, soit environ 66\,km par semaine. C'est un volume d'ultra solide,
en phase de construction au moment de l'export (pic r\'ecent \`a 130\,km \'equivalents en juin).

\subsection{Comment vos 18~mois de donn\'ees sont utilis\'es}
Toutes les sorties ne servent pas \`a la m\^eme chose. C'est le c\oe ur de la m\'ethode :
\begin{itemize}[leftmargin=1.4em]
\item \textbf{Le gros du volume} (sorties faciles et mod\'er\'ees, l'immense majorit\'e) ne renseigne pas votre
\emph{plafond} : il nourrit l'estimation de votre \textbf{endurance de base} et surtout de votre \textbf{durabilit\'e}
(votre capacit\'e \`a ne pas vous d\'egrader avec les heures).
\item \textbf{Les rares efforts quasi maximaux} (tempos, c\^otes dures, fins de course rapides) peuplent la
\og~courbe record\,\fg\ et fixent votre \textbf{vitesse critique}.
\item \textbf{Les sorties les plus longues} --- vos ultras --- jouent un r\^ole \`a part : elles sondent des dur\'ees
introuvables \`a l'entra\^inement court et servent \`a \textbf{calibrer la queue longue} du mod\`ele et la durabilit\'e.
Ce sont elles qui ancrent la pr\'ediction du 100M (partie~\ref{sec:pred}).
\end{itemize}

\subsection{Vitesse critique et courbe record}
Pour comparer des efforts men\'es sur des terrains diff\'erents, chaque seconde de course est convertie en
\textbf{vitesse \'equivalente \`a plat} via le m\^eme co\^ut de pente~\eqref{eq:minetti}. On extrait alors, pour chaque
dur\'ee, votre meilleure performance : c'est la \emph{courbe record} (figure~\ref{fig:record}). Le mod\`ele de la
\textbf{vitesse critique}~\cite{poole2016,jones2017} ajuste la relation distance--temps \(d = \VC\cdot t + D'\), o\`u
\(\VC\) est la fronti\`ere d'effort soutenable et \(D'\) une r\'eserve. Sur vos efforts plats \og~propres\,\fg\ (o\`u
l'ajustement de pente est fiable), on obtient :
\begin{equation}
\boxed{\;\VC = 2{,}91\;\text{m\,s}^{-1} \;=\; 10{,}5\kmh \;=\; 5\!:\!43\,\text{/km}\quad(\pm\,0{,}12)\;}
\end{equation}

\begin{figure}[ht]
  \centering
  \includegraphics[width=\linewidth]{figures/record.png}
  \caption{Courbe record ajust\'ee \`a la pente. En terracotta, les efforts plats propres servant \`a estimer la
  vitesse critique ; en gris, les points courts (non utilis\'es, voir encadr\'e) ; losanges verts, vos vrais ultras :
  ils se situent tr\`es en dessous de la vitesse critique, l\`a o\`u se court un 100M.}
  \label{fig:record}
\end{figure}

\begin{llattention}[Une honn\^etet\'e m\'ethodologique]
Sur sentier, les efforts \emph{courts} (moins de 30\,min) ne permettent pas d'estimer proprement la vitesse critique :
ils sont presque toujours men\'es en c\^ote, et l'ajustement de pente --- valable en r\'egime a\'erobie --- surestime
alors la vitesse \'equivalente d'un effort bref limit\'e par la puissance musculaire. Ces points (en gris) sont donc
\textbf{\'ecart\'es}. Ce n'est pas g\^enant ici : un 100M se court tr\`es loin sous la vitesse critique, o\`u ce qui
compte est l'endurance et la durabilit\'e, pas la r\'eserve ana\'erobie.
\end{llattention}

\subsection{Exposant d'endurance et durabilit\'e}
Deux param\`etres ach\`event le portrait. L'\textbf{exposant d'endurance}~\cite{riegel1981,drake2024} d\'ecrit la
vitesse de d\'eclin avec la dur\'ee ; ajust\'e sur votre enveloppe longue (30\,min \`a 6\,h), il vaut \(E \approx 1{,}22\)
--- une valeur \'elev\'ee, typique des profils d'ultra capables de tenir longtemps une fraction \'elev\'ee de leur
potentiel. La \textbf{durabilit\'e}~\cite{maunder2021,jones2024} se mesure par le \emph{d\'ecouplage} entre vitesse
ajust\'ee et fr\'equence cardiaque sur vos longues sorties : votre efficacit\'e en seconde moiti\'e baisse d'environ
\textbf{19 \`a 24\,\%} sur vos ultras de $\sim$20\,h. C'est cette d\'erive que le plan anticipe.

\begin{llnote}[Le jumeau, en r\'esum\'e]
Vitesse critique \(\VC \approx 10{,}5\kmh\) \,\textbullet\,\ r\'eserve \(D'\) faiblement contrainte (peu utile ici)
\,\textbullet\,\ exposant d'endurance \(E\approx 1{,}22\) \,\textbullet\,\ durabilit\'e : d\'erive d'efficacit\'e
\(\sim 20\,\%\) sur 20\,h \,\textbullet\,\ base : $\sim$66\,km/semaine sur 18~mois.
\end{llnote}

% =====================================================================
\section{Pr\'ediction et validation}
\label{sec:pred}

\subsection{Le principe : auto-coh\'erence dur\'ee \(\leftrightarrow\) vitesse}
Plus une course est longue, plus la vitesse soutenable baisse --- mais la dur\'ee d\'epend elle-m\^eme de cette vitesse.
On r\'esout donc un point fixe. \`A partir de vos huit ultras r\'eellement engag\'es, on ajuste la vitesse ajustée
moyenne de course \(v\) (en \kmh) en fonction de la dur\'ee \(T\) (heures) et de la densit\'e de d\'enivel\'e
\(\delta\) (m\,D+ par km) :
\begin{equation}
v(T) = 8{,}56 - 0{,}350\,\ln T - 0{,}0148\,\delta,
\qquad T = \Deq / v(T).
\label{eq:reg}
\end{equation}
Pour le 100M (\(\Deq=200{,}1\,\)km, \(\delta=53\,\)m/km), la solution auto-coh\'erente donne une vitesse ajust\'ee
moyenne de \textbf{6,58\kmh} et un temps de \textbf{30,4\,h}.

\subsection{R\'esultat}
La propagation de l'incertitude (r\'esidu du mod\`ele et al\'ea de course) par simulation de Monte-Carlo donne :
\begin{equation}
\boxed{\;T_{100M} = 30\,\text{h}\,24 \quad\text{; intervalle 80\,\% : } 29\,\text{h}\,36 - 31\,\text{h}\,18\;}
\end{equation}
soit une arriv\'ee m\'ediane le samedi vers 19\,h\,23. Rapport\'e \`a la barri\`ere de 50\,h, la marge est tr\`es
confortable : l'enjeu est la \emph{qualit\'e} de la course, pas le couperet.

\subsection{Validation crois\'ee : la preuve sur vos propres courses}
Une pr\'ediction ne vaut que si elle se v\'erifie. On applique donc une validation \emph{leave-one-out} : pour chacun
de vos huit ultras, on r\'e-estime le mod\`ele \emph{en excluant cette course}, puis on pr\'edit son temps. La
comparaison pr\'ediction/r\'ealit\'e (figure~\ref{fig:valid}) donne une erreur moyenne de \textbf{2,8\,\%}
(\'ecart-type 3,4\,\%) --- une fid\'elit\'e remarquable pour des dur\'ees de 10 \`a 21~heures.

\begin{figure}[ht]
  \centering
  \includegraphics[width=0.62\linewidth]{figures/validation.png}
  \caption{Validation crois\'ee hors \'echantillon sur vos huit ultras. Chaque point est pr\'edit par un mod\`ele
  qui ne l'a jamais \og~vu\,\fg. La quasi-totalit\'e tombe dans la bande \(\pm5\,\%\).}
  \label{fig:valid}
\end{figure}

\begin{llexemple}
Vos deux plus longues courses --- 124\,km/20\,h\,29 et 121\,km/21\,h\,15 --- sont pr\'edites hors \'echantillon \`a
\textbf{+2,5\,\%} et \textbf{+2,2\,\%}. Ce sont pr\'ecis\'ement les efforts les plus proches du 100M, ce qui renforce
la confiance dans l'extrapolation \`a 30\,h.
\end{llexemple}

% =====================================================================
\section{Le plan de pacing, segment par segment}

\subsection{Strat\'egie}
Le principe directeur est de tenir un \textbf{effort \'energ\'etique r\'egulier} (et non une allure r\'eguli\`ere) :
gr\^ace \`a l'ajustement de pente, viser une vitesse \'equivalente constante revient \`a d\'epenser au m\^eme rythme
en mont\'ee, sur le plat et en descente~\cite{swain1997}. On y superpose une \textbf{d\'erive ma\^itris\'ee} de
\(-15\,\%\) entre le d\'ebut et la fin, calibr\'ee sur votre durabilit\'e : il est normal --- et optimal --- de
ralentir l\'eg\`erement, \`a condition que ce soit \emph{pr\'evu} et non subi. L'intensit\'e moyenne ($\approx$\,63\,\%
de la vitesse critique) reste franchement dans le domaine soutenable.

Concr\`etement : \textbf{ne pas surcourir la grande mont\'ee initiale} (segments 2--3), \textbf{marcher efficacement}
les pentes raides (au-del\`a de $\sim$20\,\%, la marche active est plus \'economique que la course), \textbf{relancer
sur le plat et les faux-plats} de la seconde moiti\'e, et \textbf{prot\'eger les quadriceps} dans les longues descentes.

\subsection{Le plan}
Le tableau~\ref{tab:plan} donne, pour chaque segment, la vitesse ajust\'ee cible, l'allure r\'eelle \`a l'\oe il
(qui varie \'enorm\'ement avec la pente), le temps de mouvement, l'arr\^et pr\'evu au ravitaillement et l'heure de
passage. La figure~\ref{fig:pacing} montre la d\'erive r\'eguli\`ere de la vitesse ajust\'ee et l'allure r\'eelle
correspondante.

\begin{table}[ht]
  \centering
  \footnotesize
  \setlength{\tabcolsep}{4pt}
  \LLzebra
  \renewcommand{\arraystretch}{1.32}
  \begin{tabular}{@{}clrrrrll@{}}
    \LLheaderrow
    \LLthead{\#} & \LLthead{Ravitaillement} & \LLthead{$v$ ajust.} & \LLthead{allure} & \LLthead{t.\,mvt} & \LLthead{arr\^et} & \LLthead{arriv\'ee} & \LLthead{p\'eriode} \\
     & & \footnotesize(km/h) & \footnotesize(/km) & \footnotesize(min) & \footnotesize(min) & & \\
1 & R1 \textperiodcentered{} km\,8,1 & \LLnum{7,57} & \LLnum{7:31} & \LLnum{61} & 5 & ven. 14:05 & jour \\
2 & R2 \textperiodcentered{} km\,16,5 & \LLnum{7,50} & \LLnum{17:06} & \LLnum{144} & 5 & ven. 16:34 & jour \\
3 & R3 \textperiodcentered{} km\,28,9 & \LLnum{7,40} & \LLnum{9:50} & \LLnum{122} & 5 & ven. 18:41 & jour \\
4 & R4 \textperiodcentered{} km\,38 & \LLnum{7,33} & \LLnum{6:28} & \LLnum{59} & 5 & ven. 19:45 & \textbf{nuit} \\
5 & R5 \textperiodcentered{} km\,50,1 & \LLnum{7,27} & \LLnum{10:23} & \LLnum{126} & 15 & ven. 22:06 & \textbf{nuit} \\
6 & R6 \textperiodcentered{} km\,63,7 & \LLnum{7,17} & \LLnum{10:00} & \LLnum{136} & 5 & sam. 00:27 & \textbf{nuit} \\
7 & R7 \textperiodcentered{} km\,68,4 & \LLnum{7,10} & \LLnum{13:39} & \LLnum{64} & 5 & sam. 01:36 & \textbf{nuit} \\
8 & R8 \textperiodcentered{} km\,83,5 & \LLnum{7,02} & \LLnum{11:15} & \LLnum{170} & 5 & sam. 04:31 & \textbf{nuit} \\
9 & R9 \textperiodcentered{} km\,93,8 & \LLnum{6,91} & \LLnum{13:31} & \LLnum{139} & 15 & sam. 07:06 & \textbf{nuit} \\
10 & R10 \textperiodcentered{} km\,109,8 & \LLnum{6,82} & \LLnum{7:27} & \LLnum{119} & 5 & sam. 09:10 & jour \\
11 & R11 \textperiodcentered{} km\,121,8 & \LLnum{6,75} & \LLnum{9:30} & \LLnum{114} & 5 & sam. 11:09 & jour \\
12 & R12 \textperiodcentered{} km\,128,3 & \LLnum{6,67} & \LLnum{17:03} & \LLnum{111} & 15 & sam. 13:15 & jour \\
13 & R13 \textperiodcentered{} km\,140,4 & \LLnum{6,61} & \LLnum{7:16} & \LLnum{88} & 5 & sam. 14:48 & jour \\
14 & R14 \textperiodcentered{} km\,147,4 & \LLnum{6,56} & \LLnum{9:12} & \LLnum{64} & 5 & sam. 15:57 & jour \\
15 & R15 \textperiodcentered{} km\,156,5 & \LLnum{6,50} & \LLnum{11:57} & \LLnum{109} & 5 & sam. 17:51 & jour \\
16 & Arriv\'ee \textperiodcentered{} km\,167,2 & \LLnum{6,44} & \LLnum{8:36} & \LLnum{92} & -- & \textbf{sam. 19:23} & \textbf{nuit} \\
  \end{tabular}
  \caption{Plan de pacing. \emph{Allure} = allure r\'eelle moyenne sur le terrain (mont\'ees lentes, descentes rapides).
  \emph{t.\,mvt} = temps de course du segment. Total : 28\,h\,39 de mouvement \(+\) 1\,h\,45 d'arr\^ets \(=\) 30\,h\,24.}
  \label{tab:plan}
\end{table}

\begin{figure}[ht]
  \centering
  \includegraphics[width=\linewidth]{figures/pacing.png}
  \caption{Vitesse ajust\'ee cible (ligne, en d\'erive contr\^ol\'ee) et allure r\'eelle par segment (barres). Les
  barres les plus hautes sont les grosses mont\'ees (segments 2 et 12).}
  \label{fig:pacing}
\end{figure}

\subsection{Temps de passage, horaires et nuit}
En cumulant temps de mouvement et arr\^ets depuis le d\'epart (vendredi 13\,h\,00), on obtient les heures de passage
et leur incertitude (figure~\ref{fig:cumul}). \`A Nice, le 25~septembre, le soleil se l\`eve \`a 7\,h\,20 et se couche
\`a 19\,h\,24~\cite{noaa}. Le sc\'enario m\'edian implique \textbf{une seule nuit} : frontale du \textbf{km\,38}
(ven. 19\,h\,45) au \textbf{km\,94} (sam. 7\,h\,06). En cas de course plus lente (vers 33\,h), prévoir une seconde
frontale ou des piles de rechange pour la fin de soir\'ee du samedi.

\begin{figure}[ht]
  \centering
  \includegraphics[width=\linewidth]{figures/cumul.png}
  \caption{Temps de passage cumul\'e (m\'ediane) et intervalle \`a 80\,\%, qui s'\'elargit avec la distance. Bande
  gris\'ee : la nuit. L'incertitude reste mod\'er\'ee gr\^ace \`a la validation sur vos ultras.}
  \label{fig:cumul}
\end{figure}

\begin{llattention}[Les allures de descente sont des \emph{plafonds}]
Sur les descentes raides et techniques (segments 4, 10, 13), l'allure indiqu\'ee est celle que permet votre
\emph{m\'etabolisme} ; le \emph{terrain} (pierriers, racines, d\'enivel\'e n\'egatif sur les quadriceps) peut imposer
plus lent, surtout en fin de course. Ne cherchez pas \`a \og~tenir\,\fg\ ces allures \`a tout prix : la trace GPX ne
conna\^it pas la nature du sol.
\end{llattention}

% =====================================================================
\section{Intensit\'e, r\'eserve et durabilit\'e}

\`A 6,58\kmh\ de vitesse ajust\'ee pour une vitesse critique de 10,5\kmh, vous courez \`a $\approx$\,63\,\% de votre
fronti\`ere d'effort. C'est un point important : \`a cette intensit\'e mod\'er\'ee, la \emph{r\'eserve ana\'erobie}
(\(D'\)) n'est jamais \og~vid\'ee\,\fg\ et ne limite pas la performance --- contrairement \`a une course courte.
Les v\'eritables facteurs limitants d'un 100M sont :
\begin{itemize}[leftmargin=1.4em]
\item la \textbf{durabilit\'e} : la d\'erive d'efficacit\'e ($\sim$20\,\% sur 20\,h dans vos donn\'ees) explique
pourquoi l'allure ralentit m\^eme \`a effort constant. Le plan l'int\`egre via la d\'erive de \(-15\,\%\) ;
\item le \textbf{ravitaillement \'energ\'etique} : sur 30~heures, le d\'eficit calorique et la gestion glucidique
deviennent d\'eterminants. Un apport r\'egulier (de l'ordre de 60--90\,g de glucides par heure, \`a tester \`a
l'entra\^inement) soutient l'intensit\'e cible ;
\item la \textbf{m\'ecanique} : l'accumulation de 10\,461\,m de descente sollicite fortement les quadriceps ;
la pr\'eparation excentrique paie sur la seconde moiti\'e.
\end{itemize}

% =====================================================================
\section{Limites et conditions du jour}

Aucun mod\`ele ne capture tout. Pour un usage responsable, voici ce que ce plan \textbf{ne sait pas} :

\begin{itemize}[leftmargin=1.4em]
\item \textbf{Votre forme le jour J est inconnue.} Les donn\'ees s'arr\^etent le 22~juin, soit $\sim$13~semaines avant
la course. Le plan suppose que vous arrivez dans l'\'etat actuel, affût\'e. \textbf{Il doit \^etre recalcul\'e en
septembre} avec vos donn\'ees \`a jour --- c'est la mise \`a jour la plus importante.
\item \textbf{La technicit\'e du terrain} n'est pas dans la trace GPX : boue, pierriers, n\'ee, passages \`a la
frontale ralentissent au-del\`a du mod\`ele.
\item \textbf{La m\'et\'eo.} L'organisation annonce des \emph{temp\'eratures n\'egatives possibles \`a 2\,000\,m} ; la
nuit en altitude (segments 5--9) peut \^etre froide et venteuse, tandis que la descente vers Nice peut \^etre chaude.
Chaleur, froid et vent modifient sensiblement l'allure soutenable.
\item \textbf{Les pentes extr\^emes} : au-del\`a de \(\pm45\,\%\), le mod\`ele de co\^ut n'est plus valide, et la
course laisse place \`a la marche (m\'ecanique diff\'erente).
\item \textbf{Le calibrage des arr\^ets} (1\,h\,45 au total, r\'epartis) est une hypoth\`ese ajustable : si vous
passez plus de temps aux ravitaillements, d\'ecalez les heures d'autant (les allures de course, elles, ne changent pas).
\end{itemize}

Ces limites ne sont pas des d\'efauts cach\'es mais des bornes assum\'ees : l'incertitude annonc\'ee (intervalle \`a
80\,\%) en tient partiellement compte, et le reste rel\`eve d'une conduite prudente le jour~J.

% =====================================================================
\section{Recommandations op\'erationnelles}

\begin{enumerate}[leftmargin=1.6em]
\item \textbf{Premi\`ere mont\'ee (R2--R3) :} brider l'effort. C'est le pi\`ege classique ; chaque minute \og~gagn\'ee\,\fg\
ici se paie au double plus tard.
\item \textbf{Nuit (km\,38 \(\rightarrow\) km\,94) :} frontale + r\'eserve de piles, couche chaude pour les cr\^etes
au-dessus de 2\,000\,m.
\item \textbf{Ravitaillement :} viser 60--90\,g de glucides/heure d\`es le d\'epart ; ne pas attendre la fringale.
\item \textbf{Descentes :} prot\'eger les quadriceps t\^ot ; l'allure du tableau est un plafond, pas un objectif.
\item \textbf{Barri\`eres horaires :} la marge sur 50\,h est large ; en cas de coup de moins bien, ralentir sans
paniquer plut\^ot que de \og~cramer\,\fg\ pour rattraper un horaire.
\item \textbf{Recalculer le plan mi-septembre} avec les donn\'ees \`a jour et, si possible, la m\'et\'eo annonc\'ee.
\end{enumerate}

% =====================================================================
\section*{Limites de responsabilit\'e}
\addcontentsline{toc}{section}{Limites de responsabilit\'e}
Ce document est un prototype d'analyse \`a vis\'ee d'aide \`a la pr\'eparation sportive. Il ne constitue pas un avis
m\'edical. Les pr\'edictions reposent sur des donn\'ees pass\'ees et des mod\`eles approximatifs ; la performance r\'eelle
d\'epend de facteurs non mod\'elis\'es (sant\'e, m\'et\'eo, terrain, nutrition). Les valeurs et coefficients cit\'es
proviennent de la litt\'erature r\'ef\'erenc\'ee et de vos donn\'ees ; ils gagnent \`a \^etre revus avant tout usage
commercial.

\printbibliography

\end{document}
